\section{Limitations and Future Work}
\label{sec:limitations}

Despite its contributions, \software{} has limitations that suggest
directions for future development:

\paragraph{\LaTeX{} rendering coverage.} The browser-side preview renderer
supports a curated subset of \LaTeX{} constructs. Specialized packages
(\texttt{tikz} for diagrams, \texttt{chemfig} for chemical structures,
\texttt{forest} for linguistic trees) and custom macros are rendered as
fallback source text rather than compiled output. Full \LaTeX{} rendering in
the browser would require integrating a complete \TeX{} engine via
WebAssembly~\cite{webassembly2017} (e.g., SwiftLaTeX~\cite{swiftlatex} or a
WebAssembly-compiled TeX Live~\cite{texlive}), which
would increase the frontend bundle size by approximately 30--50 MB. A hybrid
approach---server-side full rendering with client-side caching for the
curated subset---is a promising middle ground that we plan to explore.

\paragraph{Citation database coverage.} Citation verification depends on
CrossRef, Semantic Scholar, and OpenAlex. While these databases cover a large
fraction of published research, they have gaps in coverage for non-English
publications, preprints, conference proceedings in niche venues, and gray
literature (technical reports, working papers, theses). Integrating additional
sources such as dblp (computer science bibliography), PubMed (biomedical
literature), and arXiv (preprints) would improve verification coverage.
Citation verification also does not detect fabricated-but-plausible references
that match real paper titles but are cited for non-existent claims---a more
subtle form of hallucination.

\paragraph{User studies.} We have not conducted formal controlled user studies
comparing \software{} against alternative writing workflows (Overleaf only,
ChatGPT-only, manual editing). A rigorous within-subjects or between-subjects
experiment measuring writing speed, citation accuracy, compilation success
rate, user satisfaction, and perceived control would provide quantitative
evidence for the design claims made in this paper. We consider this the most
important direction for future work and invite collaboration with HCI
researchers.

\paragraph{Collaboration support.} The current system is designed for
single-user workflows with optional token-based project sharing. Real-time
collaboration similar to Overleaf (multiple simultaneous editors, cursor
presence, operational conflict resolution) would require operational transform
or CRDT-based synchronization~\cite{crdt2011}, which introduces significant engineering
complexity. We have built a preliminary collaboration infrastructure
(WebSocket-based merge, token authorization, debounced flush), but full
real-time collaborative editing remains future work.

\paragraph{Scalability.} The filesystem-backed project model works well for
individual manuscripts but would need enhancement for institutional deployment
with hundreds of projects and users. A hybrid model---filesystem storage for
manuscript content with a metadata index in SQLite or PostgreSQL for search
and discovery---could address this while preserving local-first benefits.


\paragraph{Evaluation scope.} All evaluation was conducted by the development
team using the authors' own manuscripts and workflows. This introduces
confirmation bias and limits external validity: researchers in other
subdisciplines, working styles, or career stages may encounter friction points
that the authors did not observe. The system was tested primarily with
English-language computer science manuscripts in \LaTeX{}; performance with
humanities, social science, or natural science writing conventions---which may
rely on different packages, citation styles, or document structures---is
unknown. Future evaluation should involve researchers external to the
development team and span multiple disciplines.

\paragraph{LLM provider dependency.} \software{} delegates text generation,
review, and polishing to external LLM providers (OpenAI, Anthropic, or
locally hosted models via Ollama~\cite{ollama}). The quality and cost of AI
assistance therefore depend on the provider's pricing, rate limits, and model
capabilities, which change frequently and are outside the system's control.
While the multi-provider architecture mitigates single-vendor lock-in, it does
not insulate users from API cost escalation or service disruption. Offline
operation via local models is possible but requires non-trivial hardware and
configuration, placing it beyond the reach of many researchers.

\paragraph{Compile error repair.} When \LaTeX{} compilation fails, the system
currently reports errors and warnings in a structured format but does not
attempt automatic repair. LLM-based error diagnosis and fix suggestion is a
promising direction: the compilation log could be fed to the LLM with the
offending source lines, and the AI could suggest corrections. This aligns
directly with the document engineering philosophy---treating broken
compilation as a system-level error that the system should help resolve, not
just report.

\paragraph{Anti-AI detection.} The system includes a writing-pattern analysis
panel (Section~\ref{sec:implementation}) that examines surface-level features
such as perplexity scores, burstiness, and repetition patterns. We emphasize
three critical limitations. First, the effectiveness of AI detection tools
against modern LLMs is contested: Liang et al.~\cite{liang2024llmwriting}
documented significant rates of both false positives (flagging human-written
text as AI-generated) and false negatives (failing to detect AI-generated
text), particularly for non-native English writers. Second, this feature must
\emph{not} be used as a tool for evading academic-integrity checks or
institutional AI-use policies; doing so would undermine scholarly integrity.
Third, the writing-pattern metrics are heuristics derived from limited
baseline corpora and are discipline-dependent, language-dependent, and
model-dependent. We view this feature as strictly experimental and recommend
that users treat its output as one data point for self-review---never as proof
that text passes or fails AI detection. The interface displays explicit
warnings to this effect (Section~\ref{sec:implementation}).
